\section{Dispersion}\label{sec:dispersion}
%% FUNCTION: Shows how the EFX rate of round robin changes with the dispersion of values, all
%% other quantities at the reference cell. Round robin uses only rankings, so its allocation is the
%% same for every level of dispersion; dispersion only changes how this fixed allocation is
%% evaluated. With equal bundle sizes the EFX rate tends to one as dispersion vanishes (proven);
%% without divisibility it does not. Since this limit depends on bundle sizes, the next section
%% turns to the number of goods per agent.

Dispersion describes how far apart the values of the goods lie. When siblings divide a set of
similar books, every book is worth about the same; when they divide an estate with one house and
many small items, the values differ greatly. In our model, dispersion is the parameter $\theta$,
which we report as the Gini coefficient of the values of a single agent
(Section~\ref{sec:instances}). We vary it over its six levels, from Gini $0.06$ ($\theta=0.1$) to
Gini $0.71$ ($\theta=1.5$), and keep everything else at the reference cell: $n=4$ agents, $m=12$
goods and Spearman $0.48$ ($\rho=0.5$).

One fact shapes the whole section. Round Robin only uses the agents' rankings of the goods, and
dispersion does not change any ranking. The allocation is therefore the same for every level of
dispersion (proven, Section~\ref{sec:evaluation}); dispersion only changes how the agents
evaluate this fixed allocation. All differences in this section arise from evaluating the same
allocations, so the comparison across levels is perfectly paired.

\subsection{EFX Rate}\label{sec:dispersion-rate}

As Figure~\ref{fig:dispersion} shows, the more the values of the goods differ, the less often
Round Robin achieves EFX. In the reference cell (solid line), every one of the $20{,}000$
instances is EFX at Gini $0.06$ ($\theta=0.1$). The EFX rate falls to $0.959$ at Gini $0.14$
($\theta=0.25$), with $817$ failures, and then drops steeply to $0.574$ at Gini $0.28$
($\theta=0.5$). From there it declines more slowly: $0.375$ at Gini $0.40$ ($\theta=0.75$),
$0.293$ at Gini $0.52$ ($\theta=1$) and $0.230$ at Gini $0.71$ ($\theta=1.5$), with $15{,}407$
failures at the last level. These results are measured.

\begin{figure}[ht]
  \centering
  \includegraphics{fig42_dispersion_rate.pdf}
  \caption{EFX rate of Round Robin against the Gini coefficient of the values, for $n=4$ agents
  and Spearman $0.48$ ($\rho=0.5$). Solid line: $m=12$ goods (reference cell, $q=0$). Dashed
  line: $m=13$ goods ($q=1$). Each point rests on $R=20{,}000$ instances.}
  \label{fig:dispersion}
\end{figure}

The dashed line adds a cell with a single additional good, $m=13$, in which one agent receives
four goods and the others three. The two lines differ sharply when the values are similar and
coincide when they are not. At Gini $0.06$ ($\theta=0.1$), the EFX rate is $1.000$ with $m=12$
but only $0.499$ with $m=13$. At Gini $0.52$ ($\theta=1$), the two cells reach $0.293$ and
$0.298$, and at Gini $0.71$ ($\theta=1.5$) the cell with the additional good is even slightly
ahead, $0.254$ against $0.230$. Without divisibility, the EFX rate also declines with
dispersion, but from a much lower starting point and much more gently. Whether the number of
goods is divisible by the number of agents therefore matters only when the values of the goods
are similar. Section~\ref{sec:divisibility} examines this interaction in detail.

To check whether the decline holds beyond the reference cell, we fixed a rule before counting.
For each pair $(n,m)$ and each level of agreement, we follow the EFX rate over the six levels of
dispersion. We exclude in advance the courses that are constant by proof: those with $m=n+1$,
where the EFX rate does not depend on dispersion at all (Section~\ref{sec:bundlesize}), and those
with full agreement and $q>0$, where it is always zero (Section~\ref{sec:divisibility-identical}).
A course counts as falling if the EFX rate decreases strictly at each of the five steps. Of the
remaining $164$ courses, $144$ fall (measured). Since we knew individual cells before fixing the
rule, we treat this count as exploratory. Among the $60$ courses with $q=0$, $57$ fall; the three
exceptions occur under full agreement, where the EFX rate reaches zero before the highest level
of dispersion and stays there. Among the $104$ courses with $q>0$, $87$ fall. In $15$ of the $17$
exceptions, the rate does not fall at the first step, from Gini $0.06$ to Gini $0.14$; it stays
flat or rises by at most $0.0024$, and then falls at every later step. The remaining two
exceptions have $n=2$ agents and $m=5$ or $m=7$ goods, where the EFX rate hardly moves at all.
In summary, the EFX rate never rises noticeably with dispersion, and it falls clearly whenever
the number of goods is divisible.

\subsection{Vanishing Dispersion with Equal Bundle Sizes}\label{sec:dispersion-limit}

The value of $1.000$ at Gini $0.06$ ($\theta=0.1$) in the reference cell is no coincidence. When
all values approach each other, a comparison of bundles reduces to a comparison of their sizes.
If all bundles have the same size, removing one good from another agent's bundle always leaves
fewer goods than one's own bundle contains.

\begin{proposition}\label{prop:theta-limit}
Let $q=0$, so that every agent receives $r=m/n\geq2$ goods, and consider any allocation in which
all bundles have the same size $r$ and which does not depend on $\theta$. Then the probability
that this allocation is EFX tends to $1$ as $\theta\to0$, for every level of agreement $\rho$.
\end{proposition}

\begin{proof}
Fix an instance, that is, fixed draws $\Zig$. Since $v_i(g)=\exp(\theta\Zig)\to1$ for every good
as $\theta\to0$, we have $v_i(A_i)\to r$, $v_i(A_j)\to r$ and
$\min_{g\in A_j}v_i(g)\to1$ for every pair $i\neq j$. Hence $\aij\to r/(r-1)>1$. Because there
are finitely many pairs, there is a $\theta_0>0$, which depends on the instance, such that
$\aij>1$ for all pairs and all $\theta<\theta_0$; the allocation is then EFX. The indicator that
the allocation is EFX therefore converges to $1$ for every instance, and by dominated convergence
its expectation, the probability of EFX, converges to $1$ as well.
\end{proof}

The proposition uses neither the agreement $\rho$ nor any property of Round Robin other than equal
bundle sizes, so it also covers the balanced random allocation of Chapter~\ref{ch:random}. We
checked it in the stored cells: at Gini $0.06$ ($\theta=0.1$), all $60$ Round Robin cells with
$q=0$, across all agents, numbers of goods and levels of agreement, reach an EFX rate of exactly
$1$, and no cell with $q>0$ does. Two limits apply. First, the threshold $\theta_0$ depends on the
instance, and there is no common threshold below which every instance is EFX. For the smallest
case with $n=2$ agents and $m=4$ goods, the example of Section~\ref{sec:rrEF1} can be turned into a
proof that the EFX rate stays below $1$ for every $\theta>0$ when $\rho<1$; for larger cases we
conjecture the same but have no proof. Second, with $q>0$ the argument fails: an agent with $r$
goods compares their bundle with a bundle of $r+1$ goods, and the limit of $\aij$ is
$r/((r+1)-1)=1$ exactly. Whether EFX holds is then decided by small differences between values,
and the EFX rate does not approach $1$ (measured; we have no proof of the limit).

Figure~\ref{fig:dispersion-limit} shows what this means across numbers of goods for $n=4$ agents.
At Gini $0.06$ ($\theta=0.1$), the EFX rate jumps to exactly $1$ at $m=8$, $12$ and $16$ and lies
between $0.122$ and $0.605$ at all other numbers of goods. At Gini $0.52$ ($\theta=1$), the
pattern has vanished: the EFX rate rises gently from $0.150$ at $m=5$ to $0.409$ at $m=16$, and
the divisible cells stand out only slightly.

\begin{figure}[ht]
  \centering
  \includegraphics{fig42_dispersion_limit.pdf}
  \caption{EFX rate of Round Robin against the number of goods $m$, for $n=4$ agents and Spearman
  $0.48$ ($\rho=0.5$), at Gini $0.06$ ($\theta=0.1$) and Gini $0.52$ ($\theta=1$). Hollow
  markers: $m$ divisible by $n$ ($q=0$). Each point rests on $R=20{,}000$ instances.}
  \label{fig:dispersion-limit}
\end{figure}

\subsection{Severity of Failures}\label{sec:dispersion-severity}

Figure~\ref{fig:dispersion-survival} shows the survival curves of the reference cell, one per
level of dispersion. Because the allocation is the same at every level, the curves compare the
same allocations under different spreads of values. The curve for Gini $0.06$ ($\theta=0.1$)
stays at $1$, since no instance fails. At Gini $0.14$ ($\theta=0.25$), failures are rare and
narrow: $0.993$ of the instances reach $\astar\geq0.9$, and the conditional $5\%$ quantile is
$0.84$. With growing dispersion, the curves move down and to the left at every level of $\alpha$.
At $\alpha=0.6$, the share of instances that reach at least this level is $1.000$, $0.975$,
$0.823$, $0.651$ and $0.455$ for Gini $0.14$ to $0.71$; the conditional $5\%$ quantile falls from
$0.84$ to $0.59$, $0.37$, $0.23$ and $0.08$, and the median of $\astar$ among the instances that
are not EFX falls from $0.96$ to $0.85$, $0.73$, $0.60$ and $0.41$. Severe failures appear only at
high dispersion: the share of instances with $\astar<0.2$ is zero up to Gini $0.28$
($\theta=0.5$), $0.002$ at Gini $0.40$, $0.024$ at Gini $0.52$ and $0.163$ at Gini $0.71$
($\theta=1.5$).

\begin{figure}[ht]
  \centering
  \includegraphics{fig42_dispersion_survival.pdf}
  \caption{Share of instances with $\astar\geq\alpha$ under Round Robin, one curve per level of
  dispersion, for $n=4$ agents, $m=12$ goods and Spearman $0.48$ ($\rho=0.5$). Each curve rests
  on $R=20{,}000$ instances. The point on each curve marks the $5\%$ quantile of $\astar$ among
  the instances that are not EFX; the grey curve has no failures.}
  \label{fig:dispersion-survival}
\end{figure}

When values are similar, Round Robin therefore fails rarely, and when it fails, it misses EFX
only narrowly. When values differ greatly, it fails more often and by much more (measured). The
reason for the second part follows from the definition of $\aij$: with a large spread, a single
valuable good in another agent's bundle can outweigh one's whole bundle, and removing the least
valuable good from that bundle hardly reduces it. We state this as a conjecture.

\paragraph{Relation to literature.}
\textcite[p.~74]{bouveret2011} map ranks to utilities with different scoring functions and
recommend a quasi-indifferent one, in which every good is worth almost the same, for settings in
which the number of goods received matters most, and a lexicographic one for settings with huge
differences in value. Our dispersion axis moves continuously between these two kinds of setting,
and Proposition~\ref{prop:theta-limit} states for EFX what their quasi-indifferent case suggests:
when values hardly differ, only bundle sizes matter. The reference cell lies in the region where
the existence of an EFX allocation is open, and we assume it there
(Section~\ref{sec:existence}).

\paragraph{Recommendation.}
If the goods are worth roughly the same and their number is divisible by the number of people,
Round Robin almost always delivers EFX, and so does any split into bundles of equal size; the
more the values differ, the less Round Robin can be relied on. Because the limit depends on
whether every agent receives the same number of goods, the next section asks how the number of
goods per agent itself affects the EFX rate.
